\section{The Lodges That Differ, and the Judgment That Does Not}

The strongest honest challenge to the Church's position has never been ``the lodges are innocent''; it has been ``the lodges are not one thing.'' A judgment formed against the Grand Orient's anticlerical politics, the challenge runs, should not bind the Rotary-adjacent lodge of the English-speaking world, which requires belief in a Supreme Being, forbids religious controversy at its meetings, and plots nothing. The Church did not brush this challenge aside; she tested it, at length, in the very years the new Code was being drafted---and the most rigorous test was run by the bishops best positioned to be sympathetic.

From 1974 to 1980 the German Bishops' Conference conducted an official dialogue with the German grand lodges---bodies of the ``regular,'' Anglo-American type---including, by the participants' account, direct examination of the rituals of the three Craft degrees. Its conclusion, published in a declaration of 28 April 1980, was negative, and on doctrinal rather than political grounds. As the German conference put it in the passages most often carried into English: ``The religious conception of the Mason is relativistic: All religions are competitive attempts to explain the truth about God which, in the last analysis, is unattainable''; and, in sum, ``[i]n-depth research on the ritual and on the Masonic mentality makes it clear that it is impossible to belong to the Catholic Church and to Freemasonry at the same time.''\endnote{Declaration of the German Bishops' Conference of 28 April 1980, issued at the close of the official 1974--1980 dialogue commission; published in the German diocesan gazettes (e.g., \emph{Amtsblatt des Erzbistums K\"oln}, June 1980, pp.~102--111); Italian translation in G.~Caprile, \emph{La Civilt\`a Cattolica} 131 (1980/III) 485--502. The quotations above are taken verbatim from the documentation transmitted to the U.S. bishops with Cardinal Bernard Law's letter of 19 April 1985 (Committee on Pastoral Research and Practices; reproduction at catholicculture.org, fetched, hashed, and read 2026-07-25). Ceiling: this article inspected neither the German original nor the Civilt\`a Cattolica pages; the German conclusions are used only as reported in these identified witnesses, and the date is carried from the academic literature citing Caprile.} The finding that mattered was structural, matching what section 6 traced: not hostility, but a relativist frame built into the symbolic practice itself, such that even a ``friendly'' lodge forms its members in a conception of religious truth a Catholic cannot share. The same 1985 American documentation drew the corollary about the proposed friendly/hostile sorting: a distinction between ``favorable, neutral or hostile'' Masonry, it observed, insinuates that only membership in a hostile branch would be inadmissible---which is precisely what the 1983 declaration denies.\endnote{Cardinal Bernard Law, letter of 19 April 1985 to the U.S. bishops, with the attached reports (the German study's conclusions, W.~Whalen's study of American Masonry, and the 1985 \emph{L'Osservatore Romano} article), witness as in the preceding note. The U.S. materials confirm the same judgment for the American lodges---``the principles and basic rituals of Masonry embody a naturalistic religion active participation in which is incompatible with Christian faith and practice''---while candidly noting Freemasonry's non-unified worldwide structure.}

Rome's settled response to lodge diversity is therefore neither denial nor adjudication of particular cases from the center. The 1985 text acknowledges ``the diversity which may exist among Masonic obediences, in particular in their declared attitude towards the Church''---and states that ``the Apostolic See discerns some common principles in them which require the same evaluation by all ecclesiastical authorities.'' The 1981 and 1983 declarations lock the mechanism: no local authority may pronounce a particular lodge exempt. The analytical distinction between Anglo-American and continental Masonry thus survives---as institutional history it is simply true---but it has been ruled out as a canonical sorting principle, on the stated ground that what makes membership incompatible is common to the family.

Honesty requires marking what this settlement does and does not establish. It is the judgment of the competent doctrinal authority, twice papally approved, resting on an examined case (the German dialogue) that was as favorable to the lodges as any available. It is not a claim that every lodge member holds relativist views, which \emph{Humanum genus} n.~11 already declined to assert; nor a finding about bodies the Church has not examined; nor immune from the observation that its historical dossier is weighted toward Europe. Masonic bodies, for their part, continue to reject the Church's characterization of their principles---the UGLE's current self-description, quoted in section 2, presents the Supreme-Being clause as inclusive breadth, not indifferentism---and a fair account records that the two descriptions of the same fact have never converged. What a Catholic cannot honestly say is that the difference among lodges was never considered: it was considered, tested against the most favorable evidence, and judged not to reach the point on which the prohibition rests.
